\chapter{有限维线性代数、二次型与行列式复习}
\label{chap:appendix-linear-algebra}

\begin{chapterguide}
本附录汇总后续多元分析、微分形式和算子理论反复使用的有限维线性代数。默认读者已学过矩阵运算；重点放在不依赖坐标的对象、基变换下的公式以及有限维假设的使用位置。
\end{chapterguide}

\section{线性映射、矩阵与基变换}

设 \(V,W\) 是域 \(F\) 上有限维向量空间。映射 \(T:V\to W\) 线性是指
\[
 T(\alpha x+\beta y)=\alpha T(x)+\beta T(y).
\]
给定 \(V\) 的基 \(\mathcal B=(e_1,\ldots,e_n)\) 与 \(W\) 的基 \(\mathcal C=(f_1,\ldots,f_m)\)，矩阵 \(A=[T]_{\mathcal C\leftarrow\mathcal B}\) 的第 \(j\) 列是 \(T(e_j)\) 的 \(\mathcal C\)-坐标，因此
\[
 [T(x)]_{\mathcal C}=A[x]_{\mathcal B}.
\]

若 \(P\) 的列是 \(V\) 的新基 \(\mathcal B'\) 在旧基 \(\mathcal B\) 中的坐标，\(Q\) 的列是 \(W\) 的新基 \(\mathcal C'\) 在旧基 \(\mathcal C\) 中的坐标，则
\[
 [x]_{\mathcal B}=P[x]_{\mathcal B'},\qquad
 [y]_{\mathcal C}=Q[y]_{\mathcal C'}.
\]
从而
\[
 [T]_{\mathcal C'\leftarrow\mathcal B'}
 =Q^{-1}[T]_{\mathcal C\leftarrow\mathcal B}P.
\]
只有在 \(V=W\)，并且定义域和值域同时采用同一组新基时，公式才化为
\[
 [T]_{\mathcal B'\leftarrow\mathcal B'}
 =P^{-1}[T]_{\mathcal B\leftarrow\mathcal B}P.
\]
因此相似矩阵表示同一线性自映射在不同基下的坐标；一般线性映射的矩阵变换是左右两边分别换基。

\begin{theorem}[秩--零化度]
若 \(T:V\to W\) 线性且 \(\dim V=n\)，则
\[
 \dim\ker T+\dim\operatorname{im}T=n.
\]
\end{theorem}

\begin{proof}
取 \(\ker T\) 的基 \(e_1,\ldots,e_k\)，扩充为 \(V\) 的基。则 \(T(e_{k+1}),\ldots,T(e_n)\) 张成像空间；若其线性组合为零，对应原向量落在核中，与基的线性无关性比较可知全部系数为零。故它们是像空间的基。
\end{proof}

核与像分别衡量方程 \(T(x)=0\) 的自由度和 \(T\) 实际达到的方向。由秩--零化度立即得到：有限维自映射 \(T:V\to V\) 单射、满射、可逆三者等价。若 \(U\subseteq V\) 是子空间，则商空间 \(V/U\) 的元素是陪集 \(x+U\)，并且
\[
 \dim(V/U)=\dim V-\dim U.
\]
映射 \(T\) 在商空间上诱导的映射是否良定义，仍须验证它把同一陪集中的向量送到同一目标陪集。

\section{线性泛函、对偶基与转置映射}

对偶空间 \(V^*=\operatorname{Hom}(V,F)\)。基 \((e_i)\) 的对偶基 \((e^i)\) 由
\[
 e^i(e_j)=\delta^i_j
\]
确定；若 \(x=\sum_i x^ie_i\)，则 \(x^i=e^i(x)\)。

线性映射 \(T:V\to W\) 的转置或对偶映射
\[
 T^*:W^*\to V^*,\qquad T^*(\varphi)=\varphi\circ T
\]
满足 \((S\circ T)^*=T^*\circ S^*\)。在对偶基下，\(T^*\) 的矩阵是 \(T\) 矩阵的转置。有限维时自然映射
\[
 V\to V^{**},\qquad x\mapsto(\varphi\mapsto\varphi(x))
\]
是同构；无限维代数对偶中它一般不是满射。

若 \(U\subseteq V\)，它的湮灭子定义为
\[
 U^\circ=\{\varphi\in V^*:\varphi(u)=0\text{ 对所有 }u\in U\}.
\]
取 \(U\) 的基并扩充为 \(V\) 的基，可见
\[
 \dim U^\circ=\dim V-\dim U.
\]
同一基计算还给出
\[
 (\operatorname{im}T)^\circ=\ker T^*.
\]
再用秩--零化度便得到 \(\operatorname{rank}T^*=\operatorname{rank}T\)。这些恒等式不依赖矩阵表示；矩阵转置只是它们在对偶基下的坐标形式。

若新旧基满足 \([x]_{\mathcal B}=P[x]_{\mathcal B'}\)，则对偶坐标满足
\[
 [\varphi]_{\mathcal B'^*}=P^{\mathsf T}[\varphi]_{\mathcal B^*}.
\]
协向量与向量采用相反方向的基变换，这也是后文上、下指标变换规律的来源。这里的 \(T^*\) 是代数对偶映射；Hilbert 空间中的伴随算子还要使用内积，通常也写作 \(T^*\)，两者应由上下文区分。

\section{行列式与定向体积}

行列式是唯一满足下列条件的函数 \(\det:F^{n\times n}\to F\)：对各列分别线性，交换两列变号，且 \(\det I=1\)。Leibniz 公式为
\[
 \det A=\sum_{\sigma\in S_n}\operatorname{sgn}(\sigma)
 \prod_{i=1}^n a_{i,\sigma(i)}.
\]
\(\operatorname{sgn}(\sigma)\) 记录置换分解为换位时的奇偶性。由交替性立即得到：两列相同则行列式为零；一列是其他列的线性组合时也为零。反过来，把每一列按标准基展开，只有指标互异的项可能保留，便从交替多线性条件推出 Leibniz 公式，因而行列式的确唯一。

\begin{theorem}
\[
 \det(AB)=\det A\det B.
\]
因此 \(A\) 可逆当且仅当 \(\det A\ne0\)，且相似矩阵行列式相同。
\end{theorem}

\begin{proof}
固定 \(A\)。把 \(B\) 的列记作 \(b_1,\ldots,b_n\)，则 \(AB\) 的列是 \(Ab_1,\ldots,Ab_n\)。因此
\[
 (b_1,\ldots,b_n)\longmapsto\det(Ab_1,\ldots,Ab_n)
\]
是交替多线性函数。由行列式的唯一性，它等于某个常数乘以 \(\det B\)；令 \(B=I\)，常数就是 \(\det A\)。故 \(\det(AB)=\det A\det B\)。

若 \(A\) 可逆，则 \(1=\det A\det A^{-1}\)，所以 \(\det A\ne0\)。反之，若 \(A\) 不可逆，其列线性相关；把一列写成其他列的线性组合并使用多线性与交替性，得到 \(\det A=0\)。这证明了可逆判据。
\end{proof}

一个有序基确定实向量空间的定向；基变换矩阵行列式为正时保定向，为负时反定向。在欧氏空间中，Gram 矩阵
\[
 G=(\langle v_i,v_j\rangle)_{i,j}
\]
的行列式给出平行体体积的平方。若 \(v_i=Ae_i\)，则
\[
 \det G=\det(A^{\mathsf T}A)=(\det A)^2,
\]
故体积伸缩因子为 \(|\det A|\)，而 \(\det A\) 的符号同时记录定向。变量代换公式中出现绝对值，是因为普通体积不携带定向。

\section{特征值、正交对角化与谱定理}

若 \(T(v)=\lambda v\) 且 \(v\ne0\)，则 \(\lambda\) 是特征值。有限维矩阵的特征值满足
\[
 \det(A-\lambda I)=0.
\]
实矩阵未必有实特征值，也未必可对角化；例如平面非平凡旋转没有实特征向量，Jordan 块只有一个特征方向。

在 \(\R^n\) 中使用标准内积与欧氏范数
\[
 \langle x,y\rangle=\sum_{i=1}^n x_i y_i,\qquad
 \|x\|=\sqrt{\langle x,x\rangle}.
\]
对任意 \(x,y\) 和实数 \(t\)，有
\[
 0\le \|x-ty\|^2
 =\|x\|^2-2t\langle x,y\rangle+t^2\|y\|^2.
\]
当 \(y\ne0\) 时取 \(t=\langle x,y\rangle/\|y\|^2\)，得到 Cauchy--Schwarz 不等式
\[
 |\langle x,y\rangle|\le \|x\|\,\|y\|.
\]
等号成立当且仅当 \(x,y\) 线性相关。它进一步推出三角不等式。

若 \(U\subseteq\R^n\) 是子空间，Gram--Schmidt 过程把 \(U\) 的任一基变成标准正交基 \(u_1,\ldots,u_k\)。于是
\[
 P_Ux=\sum_{j=1}^k\langle x,u_j\rangle u_j
\]
满足 \(P_Ux\in U\) 且 \(x-P_Ux\perp U\)。分解
\[
 x=P_Ux+(x-P_Ux)
\]
唯一，因为 \(U\cap U^\perp=\{0\}\)。这一定义不依赖所选标准正交基，并给出正交投影 \(P_U\)。

\begin{theorem}[实对称谱定理]
若 \(A=A^{\mathsf T}\in\R^{n\times n}\)，则存在正交矩阵 \(Q\) 与实对角矩阵 \(\Lambda\)，使
\[
 A=Q\Lambda Q^{\mathsf T}.
\]
\end{theorem}

\begin{proof}
连续函数 \(q(x)=x^{\mathsf T}Ax\) 在单位球面上取得最大值 \(\lambda\)。这里使用有限维单位球面的紧致性；该事实在欧氏空间的紧致性理论中证明，本附录不以谱定理反过来证明它。取单位最大点 \(v\)。对任意 \(h\perp v\) 和 \(t\in\R\)，最大性给出
\[
 q(v+th)\le \lambda\|v+th\|^2.
\]
利用 \(q(v)=\lambda\)、\(\langle v,h\rangle=0\) 展开可得
\[
 2t\,h^{\mathsf T}Av+t^2
 \bigl(h^{\mathsf T}Ah-\lambda\|h\|^2\bigr)\le0.
\]
分别取正、负且趋于零的 \(t\)，一次项只能为零。因此
\[
 h^{\mathsf T}Av=0\qquad(h\perp v).
\]
所以 \(Av\in(v^\perp)^\perp=\operatorname{span}\{v\}\)，而与 \(v\) 取内积又得 \(Av=\lambda v\)。

对称性保证 \(v^\perp\) 在 \(A\) 下不变，因为
\[
 \langle Ax,v\rangle=\langle x,Av\rangle=0
 \qquad(x\perp v).
\]
在 \(v^\perp\) 上归纳，得到一组标准正交特征向量；把它们作为 \(Q\) 的列即得结论。无限维单位球面一般不紧，因此这条有限维证明不能直接用于一般无限维算子。
\end{proof}

不同特征值对应的特征向量正交，因为
\[
 \lambda\langle x,y\rangle=\langle Ax,y\rangle
 =\langle x,Ay\rangle=\mu\langle x,y\rangle.
\]

\section{二次型、正定性与 Sylvester 判据}

实二次型可写为 \(q(x)=x^{\mathsf T}Ax\)，其中只需取对称部分 \((A+A^{\mathsf T})/2\)。称 \(A\) 正定，如果 \(x^{\mathsf T}Ax>0\) 对所有 \(x\ne0\) 成立；把严格不等式改为 \(\ge0\) 得到半正定。若二次型同时取正值和负值，则称为不定。

基变换 \(x=Py\) 把矩阵变成
\[
 A\longmapsto P^{\mathsf T}AP,
\]
这称为合同变换，并不等同于相似变换。合同保持二次型的正、负、零方向数目。谱定理给出的正交对角化既是相似变换，也是合同变换。

\begin{theorem}
对实对称矩阵 \(A\)，下列条件等价：
\begin{enumerate}
 \item \(A\) 正定；
 \item \(A\) 的全部特征值为正；
 \item 存在可逆 \(R\) 使 \(A=R^{\mathsf T}R\)；
 \item 所有顺序主子式 \(\Delta_k=\det A_{1:k,1:k}\) 均为正。
\end{enumerate}
\end{theorem}

\begin{proof}
谱分解给出 \(q(x)=\sum_i\lambda_i y_i^2\)，故前两项等价。正特征值取平方根得到 \(A=R^{\mathsf T}R\)；反向由 \(\|Rx\|^2>0\) 得正定。最后一项是 Sylvester 判据：无主元交换的对称消元给出
\[
 A=LDL^{\mathsf T},\qquad
 d_k=\Delta_k/\Delta_{k-1}.
\]
下面说明这一分解为何合法。写
\[
 A=\begin{pmatrix}a&r^{\mathsf T}\\ r&C\end{pmatrix}.
\]
若首个顺序主子式 \(a=\Delta_1\ne0\)，则直接相乘可得
\[
 A=
 \begin{pmatrix}1&0\\ r/a&I\end{pmatrix}
 \begin{pmatrix}a&0\\0&C-rr^{\mathsf T}/a\end{pmatrix}
 \begin{pmatrix}1&r^{\mathsf T}/a\\0&I\end{pmatrix}.
\]
Schur 补 \(C-rr^{\mathsf T}/a\) 的第 \(j\) 个顺序主子式是
\(\Delta_{j+1}/\Delta_1\)。若所有 \(\Delta_k>0\)，可归纳继续消元，得到单位下三角矩阵 \(L\) 和
\[
 D=\operatorname{diag}(d_1,\ldots,d_n),
 \qquad d_k=\frac{\Delta_k}{\Delta_{k-1}}>0,
 \quad \Delta_0=1.
\]
于是 \(x^{\mathsf T}Ax=(L^{\mathsf T}x)^{\mathsf T}D(L^{\mathsf T}x)>0\)。

反过来，若 \(A\) 正定，则它在前 \(k\) 个坐标张成的子空间上的限制仍正定。对应矩阵 \(A_{1:k,1:k}\) 的全部特征值为正，故
\[
 \Delta_k=\det A_{1:k,1:k}>0.
\]
这证明了 Sylvester 判据，也完成四个条件的等价性。
\end{proof}

半正定情形不能只把 Sylvester 判据中的 \(>\) 改成 \(\ge\)：顺序主子式全非负并不足以控制所有方向。正确判据是全部主子式非负，或等价地全部特征值非负。退化二次型也说明二阶极值判别在 Hessian 半正定时为何不能给出严格极值结论。

\section{多线性映射与张量记号初步}

映射 \(T:V_1\times\cdots\times V_k\to W\) 若对每个变量分别线性，称为 \(k\)-线性映射。双线性型在基下写成 \(B(x,y)=B_{ij}x^iy^j\)。Einstein 求和约定把重复的一上一下指标自动求和。

张量积可以直接构造。令 \(F^{(V\times W)}\) 为以形式符号 \([v,w]\) 为基的自由向量空间，再除去由
\[
\begin{aligned}
 [v+v',w]-[v,w]-[v',w],&\qquad
 [v,w+w']-[v,w]-[v,w'],\\
 [\lambda v,w]-\lambda[v,w],&\qquad
 [v,\lambda w]-\lambda[v,w]
\end{aligned}
\]
生成的子空间。所得商空间记为 \(V\otimes W\)，\([v,w]\) 的陪集记为 \(v\otimes w\)。商掉的关系恰好保证映射
\[
 \tau:V\times W\to V\otimes W,\qquad
 \tau(v,w)=v\otimes w
\]
双线性。

若 \(B:V\times W\to U\) 双线性，先在线性基 \([v,w]\) 上规定
\(\widehat B([v,w])=B(v,w)\)。双线性保证上述关系生成的子空间包含在
\(\ker\widehat B\) 中，所以 \(\widehat B\) 唯一下降为线性映射
\[
 \widetilde B:V\otimes W\to U,\qquad
 \widetilde B(v\otimes w)=B(v,w).
\]
纯张量张成整个商空间，故这样的 \(\widetilde B\) 唯一。这就是张量积的泛性质。

若 \(e_1,\ldots,e_m\) 与 \(f_1,\ldots,f_n\) 分别是 \(V,W\) 的基，双线性展开表明 \(e_i\otimes f_j\) 张成 \(V\otimes W\)。为证线性无关，对每个 \((i,j)\) 定义双线性坐标函数
\[
 B_{ij}(v,w)=e^i(v)f^j(w).
\]
相应的 \(\widetilde B_{ij}\) 在 \(e_p\otimes f_q\) 上取值
\(\delta_{ip}\delta_{jq}\)，故任何线性关系的各个系数都为零。因此这些纯张量构成基，并且
\[
 \dim(V\otimes W)=\dim V\,\dim W.
\]
坐标变换时，每个逆变指标随基变换，每个协变指标随对偶基变换；张量记号的价值在于把这种变换律编码进指标位置。

\section*{习题}
\addcontentsline{toc}{section}{习题}

\begin{enumerate}[label=\textbf{B.\arabic*.}]
 \exerciseitem{B.1} 证明同一线性算子在两组基下的矩阵相似。
 \exerciseitem{B.2} 证明秩--零化度定理，并推出有限维空间上的单射线性自映射必为满射。
 \exerciseitem{B.3} 计算基变换下对偶基的变换矩阵。
 \exerciseitem{B.4} 证明 \((S\circ T)^*=T^*\circ S^*\) 及 \(\operatorname{rank}T^*=\operatorname{rank}T\)。
 \exerciseitem{B.5} 从交替多线性定义推出相同行的方阵行列式为零。
 \exerciseitem{B.6} 证明 \(\det(P^{-1}AP)=\det A\) 及迹的相似不变性。
 \exerciseitem{B.7} 给出不可对角化实矩阵及没有实特征值的实矩阵，并验证。
 \exerciseitem{B.8} 证明实对称矩阵不同特征值的特征空间正交。
 \exerciseitem{B.9} 证明正定矩阵的全部顺序主子式为正。
 \projectexerciseitem{B.10} \ 【前置：正交投影；预计用时：60分钟】补全实对称谱定理的极值证明。
 \optionalexerciseitem{B.11} 证明双线性型只由矩阵的对称部分决定其二次型，并用极化恒等式恢复对称双线性型。
 \optionalexerciseitem{B.12} 构造 \(V\otimes W\) 的基并证明其线性无关性。
\end{enumerate}
